\documentclass[11pt]{article}

\usepackage[T1]{fontenc}
\usepackage[utf8]{inputenc}
\usepackage{amsmath,amssymb,amsthm,mathtools}
\usepackage{booktabs}
\usepackage[margin=1in]{geometry}
% This document remains portable to minimal TeX installations that do not
% provide scalable T1 Computer Modern fonts: protrusion is harmless there,
% whereas automatic font expansion is not.
\usepackage[protrusion=true,expansion=false]{microtype}
\usepackage[hidelinks]{hyperref}

\newtheorem{theorem}{Theorem}[section]
\newtheorem{proposition}[theorem]{Proposition}
\newtheorem{lemma}[theorem]{Lemma}
\newtheorem{corollary}[theorem]{Corollary}
\theoremstyle{definition}
\newtheorem{definition}[theorem]{Definition}
\newtheorem{example}[theorem]{Example}
\theoremstyle{remark}
\newtheorem{remark}[theorem]{Remark}

\newcommand{\N}{\mathbb{N}}
\newcommand{\Q}{\mathbb{Q}}
\newcommand{\Comp}{\mathcal{C}}
\newcommand{\Chains}{\mathcal{P}}
\newcommand{\wt}{\operatorname{wt}}
\newcommand{\len}{\operatorname{len}}

\title{A Non-Recursive Composition Formula for\\
Inverse-Dyadic Values of the Fabius Function}
\author{}
\date{}

\begin{document}

\maketitle

\begin{abstract}
Let $F_n=F(2^{-n})$ denote the inverse-dyadic values of the Fabius
function.  Starting from the triangular recurrence
\[
 F_n=\frac{2^{-\binom n2}}{2^n-1}
 \sum_{k=0}^{n-1}
 \frac{2^{\binom k2}}{(n-k+1)!}F_k,
 \qquad n\geq 1,
\]
with $F_0=1$, we derive a finite, fully non-recursive solution.  After the
normalization $a_n=2^{\binom n2}F_n$, every term becomes the total weight of
all increasing paths from $0$ to $n$.  Equivalently, it is a sum over the
ordered compositions of $n$.  We prove the formula by decomposing a
composition at its last block, give equivalent increasing-chain and
nilpotent-matrix forms, and verify the values through $n=4$.  The argument is
purely finite: no limiting process or convergence theorem is involved.  We
also formalize the path and composition arguments in Lean.
\end{abstract}

\section{The recurrence and its range}

We study a sequence $(F_n)_{n\geq 0}$ with initial value
\begin{equation}\label{eq:F0}
 F_0=1
\end{equation}
and recurrence
\begin{equation}\label{eq:F-recurrence}
 F_n=\frac{2^{-\binom n2}}{2^n-1}
 \sum_{k=0}^{n-1}
 \frac{2^{\binom k2}}{(n-k+1)!}F_k,
 \qquad n\geq 1.
\end{equation}
For the Fabius function, $F_n=F(2^{-n})$, so that $F_0=F(1)=1$.
Everything below is algebraic and applies to any sequence satisfying
\eqref{eq:F0}--\eqref{eq:F-recurrence}.

The restriction $n\geq1$ is essential.  At $n=0$ the denominator $2^n-1$
vanishes and the displayed sum is empty.  Thus \eqref{eq:F-recurrence} is
not a formula for $F_0$; the initial value \eqref{eq:F0} must be supplied
separately.  In a formal field such as Lean's rational or real numbers,
division by zero is totalized, but that convention does not repair the
identity: its right-hand side would reduce to zero while $F_0=1$.

For every $n\geq1$, all quantities in \eqref{eq:F-recurrence} are defined
and its right-hand side involves only earlier values $F_0,\ldots,F_{n-1}$.
It therefore determines the sequence uniquely.  Our aim is to replace this
recursive determination by one finite multiple sum containing no $F_k$.

\section{Normalization}

The triangular powers of two in \eqref{eq:F-recurrence} are designed to
telescope.  This is clearest after a normalization.

\begin{definition}\label{def:a}
For $n\geq0$, put
\begin{equation}\label{eq:a-def}
 a_n=2^{\binom n2}F_n.
\end{equation}
\end{definition}

Since $\binom02=0$, one has $a_0=1$.  Multiplying
\eqref{eq:F-recurrence} by $2^{\binom n2}$ and replacing
$2^{\binom k2}F_k$ by $a_k$ gives
\begin{equation}\label{eq:a-recurrence}
 a_n=\frac{1}{2^n-1}
 \sum_{k=0}^{n-1}\frac{a_k}{(n-k+1)!},
 \qquad n\geq1.
\end{equation}
Equivalently, if
\begin{equation}\label{eq:Tweight}
 T_{n,k}=\frac{1}{(2^n-1)(n-k+1)!}
 \qquad (0\leq k<n),
\end{equation}
then
\begin{equation}\label{eq:a-weighted}
 a_n=\sum_{k=0}^{n-1}T_{n,k}a_k.
\end{equation}
The normalized recurrence is a weighted path recurrence on the acyclic
directed graph whose vertices are the nonnegative integers and which has an
edge $k\to n$ whenever $k<n$.

\begin{proposition}[Triangular uniqueness]\label{prop:uniqueness}
There is at most one sequence over $\Q$ or $\mathbb{R}$ satisfying
$a_0=1$ and \eqref{eq:a-recurrence}.  Consequently there is at most one
sequence satisfying \eqref{eq:F0} and \eqref{eq:F-recurrence}.
\end{proposition}

\begin{proof}
If two normalized sequences agree through index $n-1$, then
\eqref{eq:a-recurrence} expresses both values at index $n$ as the same
linear combination of those earlier values.  Induction proves equality at
every index.  The change of variables \eqref{eq:a-def} is invertible, so the
same conclusion holds for $(F_n)$.
\end{proof}

\section{The composition formula}

An \emph{ordered composition} of $n$ is a finite tuple of positive integers
whose sum is $n$.  We write
\[
 \rho=(r_1,\ldots,r_m)\vDash n
\]
and call $m=\len(\rho)$ its length.  Let $\Comp(n)$ be the set of all such
compositions.  We use the standard convention
\[
 \Comp(0)=\{()\},
\]
where $()$ is the empty composition.  For a nonempty composition define its
partial sums
\begin{equation}\label{eq:partial-sums}
 s_j=r_1+\cdots+r_j,\qquad 1\leq j\leq m.
\end{equation}
Thus $s_m=n$.  Define the weight
\begin{equation}\label{eq:composition-weight}
 \wt(\rho)=
 \prod_{j=1}^{m}
 \frac{1}{(2^{s_j}-1)(r_j+1)!}.
\end{equation}
For the empty composition the product is empty and has value $1$.

\begin{theorem}[Non-recursive composition formula]\label{thm:composition}
For every $n\geq0$,
\begin{equation}\label{eq:main-composition}
 \boxed{
 F_n=2^{-\binom n2}
 \sum_{\rho\in\Comp(n)}
 \prod_{j=1}^{\len(\rho)}
 \frac{1}{(2^{s_j}-1)(r_j+1)!}}
\end{equation}
where $\rho=(r_1,\ldots,r_m)$ and $s_j$ is given by
\eqref{eq:partial-sums}.  At $n=0$, the right-hand side is interpreted using
the unique empty composition and equals $1$.
\end{theorem}

\begin{proof}
Let
\[
 A_n=\sum_{\rho\in\Comp(n)}\wt(\rho).
\]
The empty-composition convention gives $A_0=1$.  Fix $n\geq1$.  Every
composition $\rho=(r_1,\ldots,r_m)$ of $n$ has a unique last block.  Put
\[
 k=r_1+\cdots+r_{m-1},
\]
where $k=0$ if $m=1$.  Then $0\leq k<n$, the initial tuple
$(r_1,\ldots,r_{m-1})$ is a composition of $k$, and the last block is
$r_m=n-k$.  Conversely, appending $n-k$ to any composition of $k$ produces
a composition of $n$ with penultimate partial sum $k$.

The last factor in \eqref{eq:composition-weight} is
\[
 \frac{1}{(2^n-1)(n-k+1)!}=T_{n,k}.
\]
All earlier factors form precisely the weight of the prefix composition.
Partitioning $\Comp(n)$ according to $k$ therefore gives
\begin{equation}\label{eq:A-recurrence}
 A_n=\sum_{k=0}^{n-1}
 \frac{A_k}{(2^n-1)(n-k+1)!}
 =\frac{1}{2^n-1}
 \sum_{k=0}^{n-1}\frac{A_k}{(n-k+1)!}.
\end{equation}
Thus $(A_n)$ satisfies the same recurrence and initial value as $(a_n)$.
By Proposition~\ref{prop:uniqueness}, $A_n=a_n$ for every $n$.  Finally,
$F_n=2^{-\binom n2}a_n$ by \eqref{eq:a-def}, proving
\eqref{eq:main-composition}.
\end{proof}

\begin{remark}[Finiteness]
For $n\geq1$, there are $2^{n-1}$ ordered compositions of $n$.  Hence
\eqref{eq:main-composition} is a finite sum of $2^{n-1}$ positive rational
terms.  It is fully non-recursive, although its number of terms grows
exponentially.  No convergence question is present.
\end{remark}

\begin{corollary}[Positivity]\label{cor:positive}
$F_n>0$ for every $n\geq0$.
\end{corollary}

\begin{proof}
The $n=0$ case is immediate.  If $n\geq1$, then every partial sum $s_j$ in
a composition is positive, so $2^{s_j}-1>0$, and every factorial is
positive.  Every summand in \eqref{eq:main-composition} is therefore
positive.
\end{proof}

\section{Increasing chains and an explicit multiple sum}

The composition formula can be written as a path sum without introducing
blocks.  For $m\geq1$, let $\Chains_m(n)$ denote the set of strictly
increasing chains
\begin{equation}\label{eq:chain}
 0=i_0<i_1<\cdots<i_m=n.
\end{equation}
The bijection between compositions and chains is
\begin{equation}\label{eq:chain-composition-bijection}
 r_j=i_j-i_{j-1},
 \qquad
 i_j=r_1+\cdots+r_j=s_j.
\end{equation}
Under this bijection, \eqref{eq:composition-weight} becomes
\begin{equation}\label{eq:chain-weight}
 \prod_{j=1}^{m}
 \frac{1}{(2^{i_j}-1)(i_j-i_{j-1}+1)!}.
\end{equation}

\begin{corollary}[Increasing-chain form]\label{cor:chains}
For $n\geq1$,
\begin{equation}\label{eq:chain-formula}
 \boxed{
 F_n=2^{-\binom n2}
 \sum_{m=1}^{n}
 \sum_{0=i_0<i_1<\cdots<i_m=n}
 \prod_{j=1}^{m}
 \frac{1}{(2^{i_j}-1)(i_j-i_{j-1}+1)!}.}
\end{equation}
\end{corollary}

Formula \eqref{eq:chain-formula} is the requested multiple sum.  Written
with only the internal indices displayed, its inner summation is
\[
 \sum_{1\leq i_1<\cdots<i_{m-1}\leq n-1}
 \prod_{j=1}^{m}
 \frac{1}{(2^{i_j}-1)(i_j-i_{j-1}+1)!},
 \qquad i_0=0,\quad i_m=n.
\]
When $m=1$, there are no internal indices and the inner sum consists of the
single chain $0<n$.

There is also a direct explanation for the common factor
$2^{-\binom n2}$.  Before normalization, the edge $k\to r$ in
\eqref{eq:F-recurrence} has weight
\begin{equation}\label{eq:original-edge}
 U_{r,k}=
 \frac{2^{\binom k2-\binom r2}}
 {(2^r-1)(r-k+1)!}.
\end{equation}
Along a chain \eqref{eq:chain}, the powers of two telescope:
\begin{align*}
 \prod_{j=1}^{m}
 2^{\binom{i_{j-1}}2-\binom{i_j}2}
 &=2^{\sum_{j=1}^{m}
 (\binom{i_{j-1}}2-\binom{i_j}2)}\\
 &=2^{\binom{i_0}2-\binom{i_m}2}
 =2^{-\binom n2}.
\end{align*}
The remaining edge factors are exactly \eqref{eq:chain-weight}.

\section{A nilpotent-matrix interpretation}

Fix a terminal index $N$ and consider matrices indexed by
$\{0,1,\ldots,N\}$.  Define the strictly lower-triangular matrix
\begin{equation}\label{eq:matrix-T}
 T^{[N]}_{i,j}=
 \begin{cases}
 \displaystyle
 \frac{1}{(2^i-1)(i-j+1)!},&0\leq j<i\leq N,\\[6pt]
 0,&j\geq i.
 \end{cases}
\end{equation}
Let
\[
 \mathbf{a}^{[N]}=(a_0,a_1,\ldots,a_N)^{\mathsf T},
 \qquad
 \mathbf{e}_0=(1,0,\ldots,0)^{\mathsf T}.
\]
The initial condition and \eqref{eq:a-weighted} say
\begin{equation}\label{eq:matrix-rec}
 \mathbf{a}^{[N]}=\mathbf{e}_0+T^{[N]}\mathbf{a}^{[N]}.
\end{equation}
Because $T^{[N]}$ is strictly lower triangular,
$(T^{[N]})^{N+1}=0$.  Therefore
\begin{equation}\label{eq:matrix-inverse}
 (I-T^{[N]})^{-1}
 =I+T^{[N]}+\cdots+(T^{[N]})^N,
\end{equation}
and \eqref{eq:matrix-rec} yields
\begin{equation}\label{eq:matrix-solution}
 \mathbf{a}^{[N]}=
 \left(I+T^{[N]}+\cdots+(T^{[N]})^N\right)\mathbf{e}_0.
\end{equation}

Taking the $n$th component, for any $n\leq N$, gives
\begin{equation}\label{eq:matrix-component}
 F_n=2^{-\binom n2}
 \sum_{m=0}^{n}(T^{[N]})^m_{n,0}.
\end{equation}
For $n>0$, the $m=0$ term vanishes.  Standard matrix multiplication shows
that $(T^{[N]})^m_{n,0}$ is the sum of the products of edge weights over all
length-$m$ increasing paths from $0$ to $n$.  Thus
\eqref{eq:matrix-component} is precisely Corollary~\ref{cor:chains}.

If $D=\operatorname{diag}(2^{\binom02},\ldots,2^{\binom N2})$, then the
matrix acting directly on $(F_0,\ldots,F_N)^{\mathsf T}$ is
$U=D^{-1}T^{[N]}D$.  Its entries are \eqref{eq:original-edge}; the diagonal
similarity explains the telescoping powers along paths.

\section{Low-index examples}\label{sec:examples}

We now expand the composition formula explicitly.  These checks both
illustrate the weighting rule and guard against shifts in the factorial or
the Mersenne factor.

\subsection*{$n=0$ and $n=1$}

For $n=0$, the sole composition is empty, so $a_0=1$ and $F_0=1$.
For $n=1$, the sole composition is $(1)$, with weight
\[
 \frac{1}{(2^1-1)(1+1)!}=\frac12.
\]
Since $2^{-\binom12}=1$, this gives
\[
 F_1=\frac12.
\]

\subsection*{$n=2$}

The compositions are $(2)$ and $(1,1)$.  Their weights are
\[
 \wt(2)=\frac{1}{(2^2-1)3!}=\frac1{18},
 \qquad
 \wt(1,1)=
 \frac1{(2^1-1)2!}\frac1{(2^2-1)2!}=\frac1{12}.
\]
Hence
\[
 a_2=\frac1{18}+\frac1{12}=\frac5{36},
 \qquad
 F_2=2^{-1}a_2=\frac5{72}.
\]

\subsection*{$n=3$}

For $(3)$, $(1,2)$, $(2,1)$, and $(1,1,1)$ the respective weights are
\[
 \frac1{168},\qquad
 \frac1{84},\qquad
 \frac1{252},\qquad
 \frac1{168}.
\]
Thus
\[
 a_3=\frac1{168}+\frac1{84}+\frac1{252}+\frac1{168}
 =\frac1{36},
 \qquad
 F_3=2^{-3}a_3=\frac1{288}.
\]

\subsection*{$n=4$}

There are eight compositions.  In the order
\[
 (4),\ (1,3),\ (2,2),\ (3,1),\
 (1,1,2),\ (1,2,1),\ (2,1,1),\ (1,1,1,1),
\]
their weights are
\[
 \frac1{1800},\quad \frac1{720},\quad \frac1{1620},\quad
 \frac1{5040},\quad \frac1{1080},\quad \frac1{2520},\quad
 \frac1{7560},\quad \frac1{5040}.
\]
Their sum and the resulting inverse-dyadic value are
\[
 a_4=\frac{143}{32400},
 \qquad
 F_4=2^{-6}a_4=\frac{143}{2073600}.
\]

For reference, the first five values are therefore
\begin{center}
\begin{tabular}{@{}ccl@{}}
\toprule
$n$ & $a_n=2^{\binom n2}F_n$ & $F_n$ \\
\midrule
0 & $1$ & $1$ \\
1 & $1/2$ & $1/2$ \\
2 & $5/36$ & $5/72$ \\
3 & $1/36$ & $1/288$ \\
4 & $143/32400$ & $143/2073600$ \\
\bottomrule
\end{tabular}
\end{center}

When $F_n=F(2^{-n})$, this reads
\[
 F(1)=1,\quad F(1/2)=\frac12,\quad F(1/4)=\frac5{72},
 \quad F(1/8)=\frac1{288},\quad
 F(1/16)=\frac{143}{2073600}.
\]

\section{Equivalence of the three viewpoints}

We summarize the logical equivalence, which is useful when choosing a form
for computation or formal proof.

\begin{proposition}\label{prop:equivalence}
For a sequence with $F_0=1$, the following descriptions are equivalent.
\begin{enumerate}
\item The sequence satisfies the recurrence \eqref{eq:F-recurrence} for
      every $n\geq1$.
\item It satisfies the composition formula \eqref{eq:main-composition} for
      every $n\geq0$.
\item Its normalization satisfies the finite matrix formula
      \eqref{eq:matrix-solution} for every terminal index $N$.
\end{enumerate}
\end{proposition}

\begin{proof}
The implication from the composition formula to the normalized recurrence is
the last-block decomposition \eqref{eq:A-recurrence}; undoing the
normalization gives \eqref{eq:F-recurrence}.  Conversely, the recurrence and
Proposition~\ref{prop:uniqueness} identify its solution with the composition
sum constructed in Theorem~\ref{thm:composition}.  Finally,
\eqref{eq:matrix-rec} is simply the normalized recurrence written for all
indices up to $N$, and the finite geometric identity
\eqref{eq:matrix-inverse} is equivalent to its unique solution.  Expanding
matrix powers gives the chain formula, which is carried to the composition
formula by the bijection \eqref{eq:chain-composition-bijection}.
\end{proof}

This proposition also clarifies the meaning of ``closed form'' here.  The
composition expression is not recursive: each summand depends only on the
blocks of one composition and on $n$.  The matrix inverse is likewise a
finite polynomial in a nilpotent matrix.  Both are exact finite solutions of
the triangular system.

\section{Relation to inverse-dyadic Fabius values}

The recurrence arises naturally from the half moments of Rvachev's
$\operatorname{up}$ function.  If $d_n$ denotes the $n$th half moment and
\[
 a_n=\frac{d_n}{n!},
\]
then the moment recurrence gives \eqref{eq:a-recurrence}.  The analytic
bridge to the Fabius function is
\[
 d_n=n!\,2^{\binom n2}F(2^{-n}),
\]
so the two definitions of $a_n$ agree.  Substituting this bridge into
\eqref{eq:a-recurrence} yields \eqref{eq:F-recurrence}, and
Theorem~\ref{thm:composition} therefore supplies a non-recursive exact
formula for every inverse-dyadic Fabius value.

There is no ambiguity between the usual bounded Fabius function and its
signed global extension in this statement.  Every argument $2^{-k}$ belongs
to $[0,1]$, including $2^0=1$, and the two versions agree throughout that
interval.  Thus the same finite composition formula applies to either
notation.

\section{Lean formalization}

The argument above has been formalized in Lean.  The implementation is in
\begin{verbatim}
FabiusFunction/FabiusInverseDyadicClosedForm.lean
\end{verbatim}
and is exported by the public \verb|FabiusFunction| facade.  It contains no
axioms or placeholders beyond Lean and Mathlib's standard logical
foundations.  The proof is split into a reusable path-sum layer, the explicit
composition weights, and the analytic Fabius corollaries.

\subsection*{Existing recurrence layer}

The pre-existing normalized sequence is represented by
\begin{verbatim}
fabiusRecurrenceSequence n = halfMoment n / n!
\end{verbatim}
with the positive-index recurrence exposed as
\begin{verbatim}
fabiusRecurrenceSequence_recurrence
  (n : Nat) (hn : 1 <= n) : ...
\end{verbatim}
The exact inverse-dyadic bridge and direct recurrence are represented by
declarations such as
\begin{verbatim}
fabius_inverse_two_pow_eq_recurrenceSequence
fabiusAtInverseTwoPow_recurrence_zpow
fabiusFunction_inverse_two_pow_recurrence_zpow
globalFabius_inverse_two_pow_recurrence
\end{verbatim}
The hypothesis $1\leq n$ remains explicit in source-facing recurrence
theorems.  The non-recursive formula itself is total at $n=0$, because it
uses the empty composition rather than the singular recurrence.

\subsection*{Generic weighted paths}

Mathlib provides the type \verb|Composition n| in
\verb|Mathlib.Combinatorics.Enumerative.Composition|.  Its
\verb|blocks| list records the positive increments of a path.  The formal
equivalence
\begin{verbatim}
Composition.splitLastEquiv
\end{verbatim}
identifies a composition of a positive $n$ with a proper predecessor $k<n$
and a composition of $k$.  The definitions \verb|Composition.pathWeight|
and \verb|Composition.pathSum| multiply edge weights and sum them over all
compositions.  The last-edge decomposition is the theorem
\begin{verbatim}
Composition.pathSum_eq_sum_range
\end{verbatim}
and the general solver
\begin{verbatim}
triangularRecurrence_eq_initial_mul_pathSum
\end{verbatim}
proves, by strong induction, that any semiring-valued triangular recurrence
is its initial value times the corresponding finite path sum.  A separate
\verb|triangularRecurrence_unique| theorem records uniqueness.

\subsection*{Explicit composition formula}

The theorem
\begin{verbatim}
Composition.pathWeight_eq_boundary_product
\end{verbatim}
rewrites a recursive list product chronologically, between consecutive
partial sums.
Since \verb|blocksFun| gives the blocks $r_j$ and \verb|sizeUpTo| gives the
partial sums $s_j$, the article's rational weight is implemented literally
as
\begin{verbatim}
def fabiusCompositionWeight (c : Composition n) : Rat :=
  product over j of
    1 / ((2^(c.sizeUpTo (j+1)) - 1) * (c.blocksFun j + 1)!)
\end{verbatim}
and \verb|fabiusCompositionSum n| sums this weight over all
\verb|Composition n|.  The principal exact statements are
\begin{verbatim}
fabiusCompositionSum_eq_pathSum
fabiusRecurrenceSequence_eq_sum_compositions
fabiusAtInverseTwoPow_eq_sum_compositions
fabiusAtInverseTwoPow_eq_composition_formula
fabiusAtInverseTwoPow_eq_composition_formula_by_length
fabiusFunction_inverse_two_pow_eq_sum_compositions
fabius_inverse_two_pow_eq_sum_compositions_zpow
globalFabius_inverse_two_pow_eq_sum_compositions
\end{verbatim}
The first fully expanded theorem is precisely \eqref{eq:main-composition}.
The second is the requested multiple-sum presentation: it first sums over
the length $1\leq m\leq n$, then over all compositions of $n$, retaining
only those of length $m$.  No infinite sum or matrix library is needed.
The generic real theorem holds for every bounded function satisfying the
Fabius equations; the canonical and signed-global versions follow from the
already proved analytic bridges.

\subsection*{Verification targets}

The mechanized theorem is valid for every $n:\N$, including $n=0$.
Exact rational evaluations through $n=4$ reproduce the values in
Section~\ref{sec:examples}.  The source has been checked directly and through
the public facade; the principal declarations have no \verb|sorry| and use
only the standard axioms \verb|propext|, \verb|Classical.choice|, and
\verb|Quot.sound|.  The real corollaries are derived from the rational
identity by the existing recurrence-sequence bridge, rather than by
reproving any analytic property of the Fabius function.

\section{Conclusion}

The normalization $a_n=2^{\binom n2}F_n$ turns the inverse-dyadic Fabius
recurrence into a weighted acyclic path recurrence.  Summing all paths gives
the exact composition formula \eqref{eq:main-composition}.  Its proof is the
elementary operation of removing the last block of a composition; its matrix
counterpart is the finite Neumann series of a nilpotent lower-triangular
matrix.  The resulting expression is explicit, positive, rational, and
fully non-recursive, and it incorporates the initial value at $n=0$ through
the empty-composition convention.

\end{document}
